\subsection{带算子的群}

定义2. 设 \(\Omega\) 为一个集合。一个群 G 连同 \(\Omega\) 在 G 上的一个作用，若该作用关于群律是分配的，则称 G 为以 \(\Omega\) .

为算子群的带算子群。在下文中， \(x^{\alpha}\) 将表示 \(\alpha\in\Omega\) 和 \(x\in G\) 的复合。分配性则由恒等式 \((xy)^{\alpha}=x^{\alpha}y^{\alpha}\) .

表达。在一个带算子群 G 中，每个算子定义了底层群结构的一个自同态；这些自同态有时被称为带算子群 G 的位似变换。

一个带算子群 G 称为交换的（或阿贝尔的），如果其群律是交换的。

在下文中，一个群 G 将被等同于以 \(\varnothing\) 为算子群的带算子群，这是通过赋予 G 以 \(\varnothing\) 在 G 上的唯一作用得到的。这使我们能够将群视为带算子群的特殊情形，并将我们即将陈述的关于后者的定义和结果应用于群。

例。在一个交换群 G 中，以乘法记号书写， \((xy)^{n} = x^{n}y^{n}\) 为所有 \(n \in Z\) （§2，第8号，方程（1））；因此 Z 在 G 上的作用 \(n \mapsto (x \mapsto x^{n})\) 连同群律一起，定义了 G 上带算子群的结构。

定义3. 设 G 和 G' 是以 Ω 为算子群的带算子群。从 G 到 G' 的带算子群同态是从群 G 到群 G' 的一个同态，使得

\[
f (x ^ {\alpha}) = (f (x)) ^ {\alpha}
\]

为所有 \(\alpha \in \Omega\) 以及所有 \(x \in G\) .

成立。带算子群 G 的自同态是群 G 的一个自同态，并且它与 G 的所有位似变换可交换。

由于带算子群 G 的两个位似变换不一定可交换，G 的一个位似变换一般不一定是带算子群 G 的自同态。

带算子群的恒等映射是带算子群的同态；两个带算子群同态的复合也是一个带算子群同态。一个映射是带算子群的同构，当且仅当它是带算子群的既单且满的同态，并且此时逆映射也是带算子群的同构。

更一般地，设 G（相应地 G'）是以 Ω（相应地 Ω'）为算子群的带算子群。设 φ 是从 Ω 到 Ω' 的一个映射。从 G 到 G' 的一个 φ-同态是从群 G 到群 G' 的一个同态，使得

\[
f (x ^ {\alpha}) = (f (x)) ^ {\phi (\alpha)}
\]

为所有 \(\alpha \in \Omega\) 以及所有 \(x \in G\) .

在本段的其余部分，我们将给定一个集合 \(\Omega\) 。除非另有说明，所考虑的带算子群将以 \(\Omega\) 作为算子集合。

